\section{Introduction}\label{sec:intro}

\subsection{The question}\label{sec:intro:question}

Consider a learning process that learns \emph{which inferences are valid}, and thereby learns to make valid inferences. The motivating design has six parts.
\begin{itemize}
\item It starts by imitating human inferences. These are the steps of formal or informal mathematical proofs, or of physics olympiad solutions.
\item It treats an argument as valid when it is a sequence of valid inferences.
\item It conditions the learner on \emph{coherence}: it should not have good arguments for both $P$ and $\neg P$.
\item It also conditions the learner on practical success. Occasionally the world reports a truth value.
\item On an inferentialist picture of meaning \citep{brandom1994making}, learning which inferences are valid is roughly learning what the words mean.
\item Arguments happen in \emph{contexts}. A physics problem says ``assume air pressure is $0$'', and from background knowledge one can prove that it is not. Arguments for $P$ and for $\neg P$ in such contexts should not be penalized, yet some coherence constraint is clearly needed.
\end{itemize}
Does anything of this shape work, even with unlimited compute? Does it give rise to a reasoner that derives many correct conclusions? Can one \emph{prove} that it does? The bar rises in three steps:
\begin{itemize}
\item a setup that provably works for formal mathematics;
\item one that would have worked for mathematics \emph{before} we knew how to formalize proofs;
\item a principled system for checking physics olympiad solutions.
\end{itemize}
Behind all of these sits a broader question: what does principled justification look like beyond mathematical proof?

This report answers these questions as precisely as we could. The main results are summarized below and collected in \cref{tab:intro:answers}. Each is stated as a theorem in the body, with a full proof in the appendix. Most theorems are paired with an impossibility result showing that its hypotheses cannot be dropped.

\subsection{The answers in brief}\label{sec:intro:answers}

\begin{enumerate}[label=\textbf{(\arabic*)}]
\item \textbf{Search makes soundness a worst-case property} (\cref{sec:search}).
  \begin{itemize}
  \item A reasoner searches for derivations whose steps its learned checker accepts. It is sound iff every accepted step is derivable (\cref{lem:search:stepwise}). One accepted bad step is fatal.
  \item For every distribution over steps and every $\varepsilon>0$, a single tonk-like schema with mass $<\varepsilon$ makes every formula derivable (\cref{thm:search:tonk}). Every PAC learner can be modified to stay PAC while producing outputs that trivialize the reasoner (\cref{cor:search:pac}).
  \item Majority-vote ensembles of individually coherent verifiers can derive $\bot$ forever without any member being refuted (\cref{thm:search:doctrinal}).
  \item Accuracy on human steps, which is what step-level reward models optimize, is therefore the wrong target. An experiment shows a step classifier with in-distribution AUC $0.995$ that a search procedure exploits to prove most of a set of false equations (\cref{sec:experiments}).
  \end{itemize}
\item \textbf{Cautious acceptance is the canonical sound verifier, and structure makes it affordable} (\cref{sec:caution}).
  \begin{itemize}
  \item Accepting exactly what every surviving hypothesis accepts is sound against every adaptive prover. It is also the largest sound policy (\cref{thm:caution:vs}).
  \item Its worst-case cost in human escalations is exactly a positive-elasticity dimension of the class (\cref{thm:caution:esc}).
  \item For rule schemas learned by anti-unification this cost is linear in step size. Unstructured classes cost exponentially (\cref{tab:caution:cost}).
  \item A Bayesian version is time-uniformly sound via Ville's inequality (\cref{thm:caution:ville}).
  \end{itemize}
\item \textbf{Imitation identifies rules but not norms} (\cref{sec:imitation}).
  \begin{itemize}
  \item Positive examples identify a schematic calculus exactly, at coupon-collector rates (\cref{thm:imitation:coupon}).
  \item Exact identification generalizes to derivations of any size (\cref{cor:imitation:ood}). This is a precise sense in which finite usage fixes inferential role for unbounded use.
  \item Systematic human errors are indistinguishable from rules at every rate (\cref{cor:imitation:systematic}).
  \end{itemize}
\item \textbf{Coherence is eliminative negative data, and its power is exactly Post-completeness} (\cref{sec:coherence}).
  \begin{itemize}
  \item A derivation of $\bot$ from a context certified consistent is a \emph{negative bag}: some step in it is invalid. Only bold learners ever receive this signal, with at most $\log_2(1/w^*)$ detections (\cref{thm:coherence:halving}).
  \item Every structural extension of classical consequence is classical or trivial. So one coherence datum pins down classical logic with no simplicity bias (\cref{thm:coherence:post,thm:coherence:cpc}, machine-checked).
  \item Coherence is blind among intermediate logics (\cref{thm:coherence:ipc}). In arithmetic it is Popperian, detecting $\Pi_1$ truth but not $\Sigma_1$ truth, and runs out at $\Sigma_2$ (\cref{thm:coherence:popper}).
  \item Carnap's categoricity problem is Gold's problem in the dual space of valuations. Single-conclusion data cannot fix the classical meanings, while finitely many bilateral data plus one coherence datum can (\cref{thm:coherence:carnap,thm:coherence:telltale}).
  \item What survives everything is the set of coherent uniform alternatives: a precise ``Kripkensteinian residue'' (\cref{thm:coherence:residue}).
  \end{itemize}
\item \textbf{An end-to-end theorem for formal mathematics} (\cref{sec:twotier}).
  \begin{itemize}
  \item The naive two-tier design (cautious assertion plus a bold red team) is provably unsound, because practice contains fallacies (\cref{prop:twotier:practice}).
  \item The corrected learner audits, then asserts. It assigns blame through the union of minimal conflicts, which is Reiter's duality (\cref{lem:twotier:reiter}), and states soundness relative to a refutation depth.
  \item It is sound against every adaptive prover and red team, and it converges to the target modulo an explicit residue (\cref{thm:twotier:main}).
  \item Each ingredient is necessary (\cref{tab:twotier:necessary}).
  \item For classical propositional logic and for complete decidable theories the residue is empty.
  \end{itemize}
\item \textbf{Informal mathematics before formalization} (\cref{sec:informal}).
  \begin{itemize}
  \item We model informal practice as the shadow of an unknown latent formalization. A bounded-gap conservative verifier is sound against adaptive provers. It must check links between contexts as well as steps, since Cauchy's sum theorem is an equivocation (\cref{prop:informal:equivocation}).
  \item Validity is identified exactly up to informal-validity equivalence, and no finer. Ontology is not identifiable: standard and nonstandard reals agree on the practice (\cref{thm:informal:identify,prop:informal:transfer}).
  \item Counterexample objects need at most $\log_2|\Hc|$ corrections, whereas paradoxes can need $|\Hc|-1$; the sorites is the tight case (\cref{thm:informal:objects,thm:informal:paradoxes}).
  \item The history Frege$\to$Russell$\to$Zermelo comes out as constrained repair selected by coverage of practice (\cref{thm:informal:frz}).
  \item A robust-core theorem explains why informal mathematics survived formalization (\cref{thm:informal:robust}).
  \end{itemize}
\item \textbf{Contexts and the principled checking of physics} (\cref{sec:physics}).
  \begin{itemize}
  \item Contexts denote filters of models from a declared deformation family. Suppositions are exact discharge. Idealizations may contradict the background.
  \item Eternalism is provably impossible (\cref{thm:physics:eternalism}). Coherence belongs on the root and the anchored contexts only (\cref{thm:physics:realizability}).
  \item Designating ``idealization plus full background'' as coherent forces a globally paraconsistent logic (\cref{prop:physics:misdesignation}). This is a formal version of the air-pressure worry.
  \item Exporting a conclusion from an idealization needs information of finite radius about the de-idealized family: there is no free export (\cref{thm:physics:nofree}). Context choice is an attack surface (\cref{thm:physics:stipulation}).
  \item A Structured Physics Solution checker is sound against adversarial solvers relative to a judgment layer (\cref{thm:physics:relative}). That layer, consisting of the reading of the problem and the adequacy of the top model, is ineliminable (\cref{thm:physics:judgment}).
  \item ``True in the context at hand'' becomes supertruth over admissible completions (\cref{thm:physics:superval}).
  \item The approach is worked through a EuPhO problem end to end (\cref{sec:physics:example}).
  \end{itemize}
\item \textbf{Two questions from the motivating notes.}
  \begin{itemize}
  \item A \emph{steeper simplicity penalty} on function induction with private randomness does produce the hoped-for kink at the simple model (\cref{thm:simplicity:kink}). But selection is a rate threshold. Rare valid rules are lost and cheap fallacies kept, and the threshold cannot be calibrated from imitation data (\cref{sec:simplicity}).
  \item The \emph{``philosophical completeness theorem''} (``a mode of talking is coherent iff there is something it could be talking about'') holds in a generalized weak sense, for positions, credences and context systems. It fails in the intended sense. Bounded-arity coherence constraints of the kind used in consistency-based probing do not characterize coherence (\cref{sec:existence}).
  \end{itemize}
\end{enumerate}

Three conclusions sit on top of these results (\cref{sec:philosophy}):
\begin{itemize}
\item every learning signal is \emph{one-sided};
\item principled justification is a \emph{set} of signals that jointly cover every direction of error, together with a declared residue;
\item mathematical and physical justification differ in degree, measured by the size of an explicit judgment layer, not in kind.
\end{itemize}

\begin{table}[ht]
\centering\small
\begin{tabularx}{\textwidth}{>{\raggedright\arraybackslash}p{4.2cm}>{\raggedright\arraybackslash}X}
\toprule
Success criterion & Status\\
\midrule
A setup that provably works for formal mathematics & \textbf{Yes, with explicit scope.} The two-tier learner (\cref{thm:twotier:main}) is sound and exactly convergent for classical propositional logic and complete decidable theories (\cref{cor:twotier:cpc,cor:twotier:decidable}). It is Popperian for arithmetic, with sharp impossibility results beyond $\Sigma_1$ (\cref{thm:twotier:arith}).\\
A setup that would have worked before formalization & \textbf{Yes, relative to realizability.} A bounded-gap verifier over latent formalizations is sound and identifies validity exactly up to informal-validity equivalence (\cref{sec:informal}). Open: affordable soundness when the learner must invent the formal language.\\
A principled system for checking physics olympiad solutions & \textbf{Yes, as a checker with a declared residue.} It is sound against adversarial solvers relative to the reading and the top model (\cref{thm:physics:relative}). That residue is provably ineliminable (\cref{thm:physics:judgment}). We deliberately built a checker, not a solver.\\
\bottomrule
\end{tabularx}
\caption{The three success criteria, and what this report establishes for each.}\label{tab:intro:answers}
\end{table}

\subsection{How the results were established}\label{sec:intro:method}

The work was carried out by many cooperating agent instances, coordinated by one orchestrating instance. It proceeded in four stages.
\begin{itemize}
\item Literature memos on eleven strands.
\item Seven theory threads, which developed the mathematics.
\item Adversarial verification. Every main result was attacked by two or three independent referees instructed to refute it, with brute-force computation on small cases where possible. The results were then repaired, and the repaired items were re-verified by fresh referees.
\item Drafting, followed by a fidelity check of each section against its verified source.
\end{itemize}
In the first verification round the referees found four false statements and about twenty significant gaps. Each was repaired, weakened or retracted. One theorem needed a second repair round. The full records are in \texttt{research/verification/} of the accompanying repository, and the repaired statements are the ones used here.

A core of discrete results is machine-checked in Lean~4 with Mathlib (\cref{app:lean}). It covers about fifty paper results in 14{,}600 lines with no \texttt{sorry}. Among them are:
\begin{itemize}
\item the Post-completeness and coherence theorems;
\item Carnap's problem in learning form, with the bilateral tell-tale;
\item tonk beyond the horizon, the doctrinal paradox and the halving bounds;
\item Reiter's blame duality, and the separation between paradoxes and objects;
\item bilateral completeness and the arity counterexample;
\item the rate threshold and the rate-inversion counterexample;
\item no free export from samples, and the Lipschitz certifier.
\end{itemize}
Experiments in a small symbolic toolkit illustrate the phenomena (\cref{sec:experiments}). They concern school algebra with real systematic errors, propositional natural deduction, adversarial proof search against learned verifiers, and a prototype physics checker.

We have tried to mark scope honestly:
\begin{itemize}
\item Many results are ``trivial once set up''. Their contribution is the setup, and we say so.
\item Several ingredients are classical, and we credit them. Examples are Gold, Angluin, Plotkin, Natarajan, Rivest--Sloan, KWIK, Ville and the prior--posterior-ratio martingale, Post, Glivenko, Carnap and the bilateralists, chunk-and-permeate, information-based complexity, de Finetti and Gaifman.
\item The soundness guarantees are relative to stated assumptions, chiefly realizability of the target in a structured hypothesis class and truthful designation of contexts. Computational cost is mostly not addressed.
\end{itemize}

\subsection{Reader's guide}\label{sec:intro:guide}

The report is long because it treats three domains and two side questions in full. Readers can take the following routes.
\begin{itemize}
\item \textbf{Core argument:} \cref{sec:setting,sec:search,sec:caution,sec:imitation,sec:coherence,sec:twotier}, about the formal case, then \cref{sec:philosophy}.
\item \textbf{Informal mathematics and the history of rigor:} \cref{sec:informal}, after \cref{sec:setting}.
\item \textbf{Contexts, idealization and physics checking:} \cref{sec:physics}. It is largely self-contained given \cref{sec:setting}.
\item \textbf{Learning theory of imitation and simplicity:} \cref{sec:simplicity}.
\item \textbf{Logic and the semantics of coherence:} \cref{sec:existence}.
\item \textbf{Experiments and code:} \cref{sec:experiments}.
\end{itemize}
\Cref{sec:open} collects open problems. Proofs are in the appendices, one per section.
